% Chương 5 — Thảo luận (Trúc). Nguồn số liệu: SERVER_CONTEXT/chapter4_evidence.md, model_variant_fold_metrics.md,
% TODO: đối chiếu mọi số metric với bảng Huy gửi cho Ch.4 trước khi nộp.
% reports/P5-accuracy-coursework-summary.md, P5-accuracy-failure-analysis.md (nhánh dev); mục "Giới hạn diễn giải"
% chuyển từ _moved_from_ch3.tex. Chỉ dùng số đã có trong evidence; không suy thêm chỉ số chưa được tính.
\chapter{Thảo luận}\label{ch:discussion}

Chương này đặt các kết quả của Chương~4 trở lại câu hỏi nghiên cứu ở Chương~\ref{ch:method}: chỉ với thông tin có tại thời điểm giao hàng, mô hình ước lượng rủi ro đánh giá thấp tốt đến đâu, và danh sách ưu tiên có giúp bộ phận chăm sóc khách hàng tiếp cận đúng khách hơn chọn ngẫu nhiên hay không. Mục~\ref{sec:disc-goals} so sánh kết quả với mục tiêu; mục~\ref{sec:disc-sw} trình bày điểm mạnh và hạn chế; mục~\ref{sec:disc-findings} nêu các phát hiện đáng chú ý. Mọi số liệu trong chương là kết quả trên các fold phát triển và mang tính thăm dò.

% =====================================================================
\section{So sánh với mục tiêu ban đầu}\label{sec:disc-goals}

Bảng~\ref{tab:disc-goals} đối chiếu từng mục tiêu với kết quả đạt được. Mốc so sánh cho PR-AUC là tỉ lệ dương, vì đó là PR-AUC kỳ vọng của một cách xếp hạng ngẫu nhiên; tỉ lệ dương của phân vùng huấn luyện là 10,13\% (Bảng~\ref{tab:split}). Mốc so sánh cho Recall@5\% là 5\%, mức recall kỳ vọng khi chọn ngẫu nhiên 5\% số đơn.

\begin{table}[htbp]\centering\small
\caption{Đối chiếu kết quả với mục tiêu}\label{tab:disc-goals}
\begin{tabular}{>{\raggedright\arraybackslash}p{4.1cm}>{\raggedright\arraybackslash}p{4.6cm}>{\raggedright\arraybackslash}p{2.3cm}>{\raggedright\arraybackslash}p{2.6cm}}
\toprule
Mục tiêu & Kết quả & Mốc so sánh & Đánh giá \\
\midrule
Xếp hạng tốt hơn ngẫu nhiên ở mức đơn & PR-AUC mức đơn trung bình 0,1465 (ba fold: 0,1425--0,1507) & khoảng 0,10 & Đạt, mức vượt khiêm tốn \\
Tìm đúng đơn khi chỉ xử lý được 5\% & Recall@5\% trung bình 0,1212 & 0,05 & Đạt, gấp khoảng 2,4 lần \\
Mục tiêu dừng của kế hoạch cải thiện & PR-AUC lăn 0,1465 & 0,25 & Chưa đạt \\
Cải thiện so với mô hình ban đầu & PR-AUC lăn tăng từ 0,1363 lên 0,1465 & mô hình ban đầu & Đạt (+7,5\%) \\
Ước lượng cuối cùng trên tập kiểm tra niêm phong & Chưa thực hiện & -- & Chưa đạt \\
Hệ thống xử lý case dựa trên điểm số & Chạy đủ luồng từ điểm số đến xử lý case & -- & Đạt, dùng điểm minh họa \\
\bottomrule
\end{tabular}
\end{table}

Kết quả trả lời câu hỏi nghiên cứu theo hai chiều. Về khả năng xếp hạng, câu trả lời là \emph{có, nhưng ở mức hạn chế}: PR-AUC cao hơn mốc ngẫu nhiên khoảng 45\%, và mức vượt này ổn định qua cả ba fold chứ không chỉ xuất hiện ở một giai đoạn. Về giá trị vận hành, Recall@5\% bằng 0,1212 nghĩa là nếu bộ phận chăm sóc khách hàng chỉ liên hệ được 5\% số đơn giao đúng hạn, danh sách theo mô hình chứa khoảng 12\% số đơn sẽ bị đánh giá thấp, gấp khoảng 2,4 lần so với chọn ngẫu nhiên. Nói cách khác, cùng một ngân sách liên hệ, mô hình giúp tiếp cận số khách có nguy cơ nhiều gấp đôi; nhưng gần 88\% khách sẽ đánh giá thấp vẫn nằm ngoài danh sách.

Mục tiêu PR-AUC 0,25 không đạt. Kế hoạch cải thiện đã thử mười giai đoạn có kiểm soát, giữ lại ba nhóm đặc trưng và dừng có chủ đích khi các thử nghiệm tiếp theo không còn cải thiện ổn định. Việc dừng ở đây là một kết quả, không phải một thiếu sót về nỗ lực: khi mô hình phức tạp hơn và đặc trưng phong phú hơn đều không vượt được mô hình đã chọn, khả năng cao là giới hạn nằm ở lượng tín hiệu có trong dữ liệu tại thời điểm giao hàng chứ không nằm ở thuật toán (mục~\ref{sec:disc-findings}).

% =====================================================================
\section{Điểm mạnh và hạn chế}\label{sec:disc-sw}

\subsection{Điểm mạnh}

\textbf{Đánh giá trong đúng điều kiện sử dụng.} Ba quyết định thiết kế làm cho con số khiêm tốn ở trên đáng tin hơn những con số cao hơn thường gặp: chia dữ liệu theo thời gian thay vì ngẫu nhiên, đặc trưng lịch sử chỉ dùng các đơn đã giao \emph{và} đã được đánh giá trước thời điểm dự đoán, và mỗi phân vùng dữ liệu chỉ phục vụ một quyết định. Các nghiên cứu chia ngẫu nhiên trên cùng bộ dữ liệu cho phép mô hình học từ tương lai~\cite{orman2026}, nên kết quả của họ và của đề tài không so sánh trực tiếp được; điều đề tài đo được là hiệu năng khi mô hình thực sự được dùng sau khi đơn vừa giao.

\textbf{Mỗi lựa chọn đều có mốc đối chứng.} Mô hình được chọn bằng trung bình ba fold rolling-origin thay vì một tập kiểm định, mỗi nhóm đặc trưng chỉ được giữ khi cải thiện chỉ số chính, và việc bỏ từng nhóm (leave-one-group-out) xác nhận cả ba nhóm được giữ đều cần thiết. Nhờ vậy, mỗi thành phần trong mô hình cuối cùng đều có bằng chứng riêng cho sự có mặt của nó.

\textbf{Kết nối được với vận hành.} Ngưỡng theo phần trăm biến năng lực xử lý thành tham số của mô hình; luật max-v1 được đánh số phiên bản; điểm số được ghi kèm phiên bản mô hình và ngưỡng; ứng dụng không cho phép sửa điểm và ghi lại kết quả xử lý có cấu trúc (mục~\ref{sec:demo}). Đây là nền tảng cần thiết để đánh giá hiệu quả can thiệp ở giai đoạn sau, điều mà phần lớn nghiên cứu dự báo sự hài lòng chưa đặt ra.

\subsection{Hạn chế}

\textbf{Kết quả chưa phải ước lượng cuối cùng.} Giai đoạn kiểm tra lịch sử đã bị lộ trong quá trình phát triển, nên mọi kết quả trên giai đoạn này được xếp loại thăm dò. Kết luận xác nhận cần một tập kiểm tra mới, được khóa trước khi xem. Cùng với đó, khoảng tin cậy bootstrap theo cụm đơn hàng, đường cong precision--recall và đường cong hiệu chỉnh trên tập kiểm tra chưa được tính.

\textbf{So sánh chưa đầy đủ.} Kho lưu trữ chỉ giữ số liệu theo fold của rừng ngẫu nhiên và ba cấu hình CatBoost; số liệu của mô hình tỉ lệ nền, hồi quy logistic và HistGradientBoosting không còn, nên chưa có một bảng so sánh công bằng giữa năm họ mô hình trên cùng bộ đặc trưng. Đề tài cũng chưa so với chính sách xếp theo giá trị đơn, một mốc đơn giản mà nghiên cứu gần đây cho thấy đôi khi khó vượt qua~\cite{li2026valuesorting}.

\textbf{Nhãn yếu và tự chọn.} Mỗi item thừa hưởng nhãn của đơn, nên trong đơn nhiều sản phẩm, mô hình học cả những item không liên quan tới sự bất mãn. Nhãn chỉ có ở các đơn có đánh giá, mà việc viết đánh giá là một quá trình tự chọn~\cite{li2008}.

\textbf{Chưa có giải thích mô hình và chưa có điểm thật trong ứng dụng.} Chưa có lần tính SHAP nào, nên chưa biết mô hình dựa vào đặc trưng nào cho từng đơn. Ứng dụng đang chạy với điểm minh họa và ngưỡng cố định (mục~\ref{sec:demo-limits}).

\subsection{Giới hạn khi diễn giải kết quả}

Kết quả của đề tài cần được hiểu trong các giới hạn sau:
\begin{itemize}
  \item Không có đánh giá không có nghĩa là khách hài lòng, do quá trình tự chọn khi viết đánh giá~\cite{li2008}.
  \item Đánh giá thấp không đồng nghĩa với việc khách rời bỏ hoặc gây thiệt hại doanh thu.
  \item Điểm rủi ro cao không có nghĩa là khách sẽ phản hồi tốt với can thiệp. Thực nghiệm cho thấy nhóm khách rủi ro cao nhất và nhóm nhạy cảm nhất với can thiệp có thể chỉ trùng nhau ở mức gần như ngẫu nhiên~\cite{ascarza2018}.
  \item Các mối liên hệ mô hình học được là liên hệ dự báo, không phải quan hệ nhân quả: một đặc trưng có ích cho dự báo không có nghĩa là thay đổi đặc trưng đó sẽ thay đổi sự hài lòng.
  \item Kết luận chỉ áp dụng cho giai đoạn dữ liệu Olist 2016--2018, trừ khi được kiểm định trên dữ liệu bên ngoài.
\end{itemize}

% =====================================================================
\section{Phát hiện đáng chú ý}\label{sec:disc-findings}

\subsection{Cách người bán vận hành quan trọng hơn điểm đánh giá trước đó}

Trong mười giai đoạn cải thiện, các nhóm đặc trưng được giữ lại là tốc độ và độ biến động vận hành của người bán, cấu trúc đơn hàng, và hình thái cùng độ đầy đủ nội dung của sản phẩm. Khi bỏ từng nhóm, PR-AUC lăn giảm lần lượt 0,0042, 0,0050 và 0,0020. Ngược lại, nhóm lịch sử mục tiêu, tức điểm đánh giá trước đó của người bán và danh mục được mã hóa mạnh hơn, chỉ cải thiện tốt nhất 0,0015 và bị loại; các nhóm địa lý, thanh toán, lịch và chất lượng dữ liệu cũng bị loại.

Kết quả này đáng chú ý vì trực giác thông thường cho rằng người bán từng bị chê sẽ tiếp tục bị chê. Dữ liệu cho thấy, khi đã có thông tin cơ bản về lịch sử đánh giá, thêm các phép mã hóa phức tạp hơn của cùng thông tin đó không giúp ích; điều giúp ích là \emph{cách} người bán vận hành, như thời gian từ duyệt đơn tới bàn giao vận chuyển và độ thất thường của thời gian giao. Một cách giải thích hợp lý là: với đơn đã giao đúng hạn, khách không phàn nàn về việc trễ, nhưng một người bán vận hành thất thường thường cũng chuẩn bị hàng và đóng gói kém ổn định. Đây là giả thuyết cần kiểm chứng bằng nội dung đánh giá, không phải kết luận nhân quả. Nhóm cấu trúc đơn được giữ lại cũng phù hợp với bằng chứng rằng giao hàng tách thành nhiều kiện làm giảm sự hài lòng~\cite{brylla2022}.

\subsection{Một tập kiểm định đơn lẻ có thể chọn sai mô hình}

Cấu hình CatBoost tốt nhất đạt PR-AUC 0,2239 trên phân vùng lựa chọn mô hình, cao hơn rõ rệt so với 0,2108 của rừng ngẫu nhiên. Nếu chọn theo một tập kiểm định duy nhất, CatBoost sẽ thắng. Nhưng trung bình trên ba fold rolling-origin, CatBoost đạt 0,1458, thấp hơn 0,1465 của rừng ngẫu nhiên, nên theo quy tắc đã khai báo trước, rừng ngẫu nhiên được giữ. Khoảng cách giữa hai cách đánh giá lớn hơn nhiều so với khoảng cách giữa hai mô hình, cho thấy với lớp dương hiếm và tỉ lệ dương dịch chuyển theo thời gian, thứ hạng của các mô hình trên một giai đoạn đơn lẻ chủ yếu phản ánh nhiễu của giai đoạn đó. Đây là bằng chứng thực nghiệm cho lựa chọn đánh giá trên nhiều fold ở Chương~\ref{ch:method}.

Cũng cần lưu ý PR-AUC trên phân vùng lựa chọn mô hình (khoảng 0,21) cao hơn trung bình ba fold (khoảng 0,15). Một phần chênh lệch đến từ tỉ lệ dương của phân vùng này cao hơn (11,92\% so với 10,13\%), vì mốc ngẫu nhiên của PR-AUC chính là tỉ lệ dương. Hai con số vì vậy không nên được đặt cạnh nhau như thể đo cùng một thứ.

\subsection{Trọng số lớp làm xác suất kém tin cậy hơn dự đoán hằng số}

Brier score trung bình ba fold của rừng ngẫu nhiên là 0,1907. Một dự đoán hằng số bằng tỉ lệ dương khoảng 0,10 cho Brier score xấp xỉ $0{,}10 \times 0{,}90 = 0{,}09$, tức tốt hơn gần một nửa. Mô hình xếp hạng tốt hơn ngẫu nhiên nhưng xác suất của nó kém tin cậy hơn một hằng số. Nguyên nhân là trọng số lớp cân bằng đẩy xác suất dự đoán lên một cách có hệ thống, đúng như cảnh báo của~\textcite{vandengoorbergh2022}; ba cấu hình CatBoost có Brier trung bình fold từ khoảng 0,13 đến 0,17 cũng cho thấy mức độ lệch phụ thuộc vào mô hình. Phát hiện này giải thích vì sao bước hiệu chỉnh xác suất ở mục~\ref{sec:models} là bắt buộc chứ không phải tùy chọn: nếu đưa xác suất chưa hiệu chỉnh lên giao diện, quản lý sẽ ước lượng số đơn có vấn đề cao hơn thực tế một cách có hệ thống. Với ứng viên CatBoost, Brier score sau hiệu chỉnh sigmoid trên phân vùng so sánh là 0,1071; số liệu tương ứng của rừng ngẫu nhiên sau hiệu chỉnh chưa được lưu trong kho và cần được bổ sung.

\subsection{Ngưỡng là quyết định vận hành, không phải tham số kỹ thuật}

Khi đối chiếu bản demo với kết quả nghiên cứu (Bảng~\ref{tab:demo-threshold}), cùng một từ ``HIGH'' mang hai nghĩa rất khác nhau. Trong bản demo với ngưỡng cố định 0,7, mức HIGH chiếm 29\% số case; theo chính sách nghiên cứu, mức HIGH là 5\% đơn có điểm cao nhất, tương ứng ngưỡng khoảng 0,286 với rừng ngẫu nhiên. Nếu mang ngưỡng 0,7 sang mô hình thật, số đơn vượt ngưỡng sẽ ít hơn nhiều so với 5\%, vì ngưỡng của nhóm 5\% cao nhất chỉ khoảng 0,286; nếu mang tỉ lệ 29\% sang vận hành, quản lý phải tự phân công gần một phần ba số đơn. Cả hai trường hợp đều làm hỏng mục đích của hàng đợi ưu tiên. Bài học rút ra là ngưỡng phải được chọn lại mỗi khi mô hình hoặc năng lực xử lý thay đổi, và nên được biểu diễn bằng tỉ lệ phần trăm số đơn thay vì một con số xác suất cố định~\cite{gomezvargas2025}.
